On dispose d'une solution aqueuse $\mathrm{(S_{A})}$ d'acide méthanoïque
$\ce{HCOOH_{(aq)}}$ de volume $V=\qty{1}{\liter}$, de concentration molaire
$C_{A}=\qty{0,10}{\M}$ et de $\mathrm{pH}=2,4$.
\begin{enumerate}
    \item Définir un acide selon Brönsted.
    \item Écrire l'équation modélisant la transformation chimique entre l'acide
          méthanoïque et l'eau.
    \item Recopier sur votre copie le tableau d'avancement et le compléter.
          \begin{table}[H]
          \renewcommand{\arraystretch}{1.5}
          \begin{tabularx}{\textwidth}{|
              >{\centering\arraybackslash}m{2.5cm}|
              >{\centering\arraybackslash}m{3.5cm}|C|C|C|C|}
              \hline
              \multicolumn{2}{|c|}{Équation chimique} &
              \multicolumn{4}{c|}{\dotfill} \\
              \hline
              État du système & Avancement de la réaction en (\unit{\mol}) &
              \multicolumn{4}{c|}{Quantité de matière en (\unit{\mol})} \\
              \hline
              État initial & 0 & \dotfill & \dotfill & \dotfill & \dotfill \\
              \hline
              État intermédiaire & $x$ & \dotfill & \dotfill & \dotfill &
              \dotfill \\
              \hline
              État final & $x_{f}$ & \dotfill & \dotfill & \dotfill &
              \dotfill \\
              \hline
          \end{tabularx}
          \end{table}
    \item Calculer la valeur de l'avancement final $x_{f}$ de cette réaction.
    \item Calculer le taux d'avancement final $\tau$ de cette réaction.
          Conclure.
    \item Montrer que le quotient de réaction à l'état d'équilibre du système
          chimique s'écrit~:
          $Q_{r,\eq}=\frac{10^{-2\cdot\mathrm{pH}}}{C_{A}-10^{-\mathrm{pH}}}$.
          Calculer sa valeur.
    \item Déduire la valeur de la constante d'équilibre $\mathrm{K}$ associé à
          l'équation de la réaction. 
\end{enumerate}


